% Chapter 4, Section 1 _Linear Algebra_ Jim Hefferon
%  http://joshua.smcvt.edu/linearalgebra
%  2001-Jun-12
\section{幂零性}
本章将证明，每个方阵都相似于一个能写成两类简单矩阵之和的矩阵。
前一小节关注了第一类简单矩阵，即对角矩阵。
现在我们来考虑另一类。









\subsectionoptional{自复合}\index{self composition!of maps}\index{map!self composition}\index{transformation!composed with itself}\index{composition!self}
% \noindent\textit{This subsection is optional, although it is necessary
% for later material in this section and in the next one.}

由于线性变换 $\map{t}{V}{V}$ 的定义域与陪域相同，我们可以将
$t$ 与自身复合：
\( t^2=\composed{t}{t} \)，
以及 \( t^3=\composed{t}{\composed{t}{t}} \)，
等等。\appendrefs{function iteration}\spacefactor=1000 %
\begin{center}
  \includegraphics{jc/mp/ch5.1}
\end{center}
对变换的迭代使用上标幂记号 $t^j$，
与我们对其方阵表示所用的记号是相容的，
因为若 $\rep{t}{B,B}=T$，则 \( \rep{t^j}{B,B}=T^j \)。


\begin{example} \label{ex:DerivIter}
对于求导映射 \( \map{d/dx}{\polyspace_3}{\polyspace_3} \)，它由下式给出
\begin{equation*}
  a+bx+cx^2+dx^3\xmapsunder{d/dx} b+2cx+3dx^2
\end{equation*}
其二次幂就是二阶导数，
\begin{equation*}
  a+bx+cx^2+dx^3\xmapsunder{d^2/dx^2} 2c+6dx
\end{equation*}
其三次幂就是三阶导数，
\begin{equation*}
  a+bx+cx^2+dx^3\xmapsunder{d^3/dx^3} 6d
\end{equation*}
而任意更高次幂都是零映射。
\end{example}

\begin{example}
空间 $\matspace_{\nbyn{2}}$（即 $\nbyn{2}$ 矩阵构成的空间）中的这个变换
\begin{equation*}
  \begin{mat}
    a  &b  \\
    c  &d
  \end{mat}
  \mapsunder{t}
  \begin{mat}
    b  &a  \\
    d  &0
  \end{mat}
\end{equation*}
有如下二次幂
\begin{equation*}
  \begin{mat}
    a  &b  \\
    c  &d
  \end{mat}
  \mapsunder{t^2}
  \begin{mat}
    a  &b  \\
    0  &0
  \end{mat}
\end{equation*}
以及如下三次幂。
\begin{equation*}
  \begin{mat}
    a  &b  \\
    c  &d
  \end{mat}
  \mapsunder{t^3}
  \begin{mat}
    b  &a  \\
    0  &0
  \end{mat}
\end{equation*}
此后，$t^4=t^2$，$t^5=t^3$，依此类推。
\end{example}

\begin{example}
考虑平移变换 $\map{t}{\C^3}{\C^3}$。
\begin{equation*}
  \colvec{x \\ y \\ z} \mapsunder{t} \colvec{0 \\ x \\ y}
\end{equation*}
我们有
\begin{equation*}
  \colvec{x \\ y \\ z} \mapsunder{t} \colvec{0 \\ x \\ y}
                       \mapsunder{t} \colvec{0 \\ 0 \\ x}
                       \mapsunder{t} \colvec{0 \\ 0 \\ 0}
\end{equation*}
所以值域空间下降到平凡子空间。
\begin{equation*}
  \rangespace{t}=\set{\colvec{0 \\ a \\ b}\suchthat a,b\in\C}
  \qquad
  \rangespace{t^2}=\set{\colvec{0 \\ 0 \\ c}\suchthat c\in\C}
  \qquad
  \rangespace{t^3}=\set{\colvec{0 \\ 0 \\ 0}}
\end{equation*}
\end{example}

这些例子表明，经过若干次迭代之后，映射就趋于稳定。

\begin{lemma}  \label{le:RangeAndNullChains}
对任意变换 \( \map{t}{V}{V} \)，各次幂的值域空间构成一个递降链
\begin{equation*}
  V\supseteq \rangespace{t}\supseteq\rangespace{t^2}\supseteq\cdots
\end{equation*}
而零空间构成一个递升链。
\begin{equation*}
  \set{\vec{0}}\subseteq\nullspace{t}\subseteq\nullspace{t^2}\subseteq\cdots
\end{equation*}
此外，存在 \( k>0 \)，使得当幂小于 $k$ 时各包含关系都是真包含：~若
$j<k$，则 $\rangespace{t^j}\supset\rangespace{t^{j+1}}$
且
$\nullspace{t^j}\subset\nullspace{t^{j+1}}$；
而当
$j\geq k$ 时，则 $\rangespace{t^j}=\rangespace{t^{j+1}}$
且
$\nullspace{t^j}=\nullspace{t^{j+1}}$）。
\end{lemma}

\noindent （$k=1$ 的情形可能发生，例如当 $t$ 是恒等映射时，
此时两条链中没有一个包含关系是真包含。）

\begin{proof}
首先回忆，对任意映射，其值域空间的维数加上其零空间的维数等于
其定义域的维数。
于是，如果值域空间的维数收缩，那么零空间的维数必定上升。
这里我们只做值域空间这一半，其余部分留作
\nearbyexercise{exer:RangeAndNullChains}。

我们先证明
值域空间构成一条链。
若 $\vec{w}\in\rangespace{t^{j+1}}$，即有某个 $\vec{v}$ 使得
$\vec{w}=t^{j+1}(\vec{v})$，
则 $\vec{w}=t^{j}(\,t(\vec{v})\,)$。
于是 $\vec{w}\in\rangespace{t^{j}}$。

接下来我们验证“进一步”的性质：
链中的包含关系起初都是真包含，
而从某个幂 $k$ 开始值域空间彼此相等。
我们先证明：如果链中任意一对相邻的值域空间相等
\( \rangespace{t^{k}}=\rangespace{t^{k+1}} \)，
那么其后所有的值域空间也都相等：
\( \rangespace{t^{k+1}}=\rangespace{t^{k+2}} \)，等等。
这是因为
\( \map{t}{\rangespace{t^{k+1}}}{\rangespace{t^{k+2}}} \)
与
\( \map{t}{\rangespace{t^{k}}}{\rangespace{t^{k+1}}} \) 是定义域相同的同一个映射，
因此它有相同的值域
\( \rangespace{t^{k+1}}=\rangespace{t^{k+2}} \)
（由归纳法可知这对所有更高的幂都成立）。
因此，值域空间的链一旦停止严格递减，
从那一点起它就稳定了。

最后我们证明这条链必定最终停止递减。
每个值域空间都是前一个的子空间。
若要成为真子空间，它的维数就必须严格更低
（见 \nearbyexercise{exer:PropSubspStrictLowerDimen}）。
这些空间都是有限维的，所以这条链只能下降有限多步。
也就是说，幂 $k$ 至多是 $V$ 的维数。
\end{proof}

\begin{example}
$\polyspace_3$ 上的求导映射 $a+bx+cx^2+dx^3\xmapsunder{d/dx} b+2cx+3dx^2$
有如下值域空间链。
\begin{equation*}
  \rangespace{t^0}=\polyspace_3
    \:\supset\:\rangespace{t^1}=\polyspace_2
    \:\supset\:\rangespace{t^2}=\polyspace_1
    \:\supset\:\rangespace{t^3}=\polyspace_0
    \:\supset\:\rangespace{t^4}=\set{\zero}
    % =\rangespace{t^5}=\set{\zero}=
    % \,=\,\cdots
\end{equation*}
链中其后的各项都是平凡空间。
它有如下零空间链。
\begin{equation*}
  \nullspace{t^0}=\set{\zero}
  \:\subset\: \nullspace{t^1}=\polyspace_0
  \:\subset\:\nullspace{t^2}=\polyspace_1
  \:\subset\:\nullspace{t^3}=\polyspace_2
  \:\subset\:\nullspace{t^4}=\polyspace_3
\end{equation*}
其后的各项都是整个空间。
\end{example}

\begin{example} \label{exam:PolyRankFalls}
设 \( \map{t}{\polyspace_2}{\polyspace_2} \) 是映射
\( d_0+d_1x+d_2x^2 \mapsto 2d_0+d_2x. \)
如引理所描述的，在迭代下值域空间收缩
\begin{equation*}
  \rangespace{t^0}=\polyspace_2
    \quad
  \rangespace{t}=\set{a_0+a_1x\suchthat a_0,a_1\in\C}
    \quad
  \rangespace{t^2}=\set{a_0\suchthat a_0\in\C}
\end{equation*}
然后稳定下来，使得 $\rangespace{t^2}=\rangespace{t^3}=\cdots$。
零空间增长
\begin{equation*}
  \nullspace{t^0}=\set{0}
    \quad
  \nullspace{t}=\set{b_1x\suchthat b_1\in\C}
    \quad
  \nullspace{t^2}=\set{b_1x+b_2x^2\suchthat b_1,b_2\in\C}
\end{equation*}
然后稳定下来，$\nullspace{t^2}=\nullspace{t^3}=\cdots$。
\end{example}

\begin{example}
投影到前两个坐标的变换 \( \map{\pi}{\C^3}{\C^3} \)
\begin{equation*}
   \colvec{c_1 \\ c_2 \\ c_3}
     \mapsunder{\pi}
   \colvec{c_1 \\ c_2 \\ 0}
\end{equation*}
满足 \( \C^3\supset\rangespace{\pi}=\rangespace{\pi^2}=\cdots \)
且 \( \set{\zero}\subset\nullspace{\pi}=\nullspace{\pi^2}=\cdots\, \)，
其中的值域空间与零空间为
\begin{equation*}
  \rangespace{\pi}=\set{\colvec{a \\ b \\ 0}\suchthat a,b\in\C}
  \qquad
  \nullspace{\pi}=\set{\colvec{0 \\ 0 \\ c}\suchthat c\in\C}
\end{equation*}
\end{example}

\begin{definition}
设 \( t \) 是 \( n \) 维空间上的一个变换。
\definend{广义值域空间}\index{generalized range space}%
\index{range space!generalized}
（或 \definend{值域空间的闭包\/}\index{closure!of range space}%
\index{range space!closure of}）
是 $\genrangespace{t}=\rangespace{t^n}$。
\definend{广义零空间}\index{generalized null space}%
\index{null space!generalized}
（或 \definend{零空间的闭包\/}\index{closure!of null space}%
\index{null space!closure of}）
是 $\gennullspace{t}=\nullspace{t^n}$。
\end{definition}

这幅图给出了说明。
横轴给出变换的幂~$j$。
纵轴把
$t^j$ 的值域空间的维数表示为零点之上的距离，
因而也同时显示了零空间的维数，
因为两者之和等于定义域的维数 $n$。
\begin{center}
  % \includegraphics{jc/mp/ch5.2}
  \includegraphics{jc/mp/ch5.8}
\end{center}
在迭代下
秩下降而零度上升，
直到存在某个 $k$ 使得
映射达到稳定状态
$\rangespace{t^k}=\rangespace{t^{k+1}}=\genrangespace{t}$
且 $\nullspace{t^k}=\nullspace{t^{k+1}}=\gennullspace{t}$。
这最迟必定在第 $n$ 次迭代时发生。
% The steady state's distance above zero is the dimension of the
% generalized range space
% and its distance below $n$ is the dimension of
% the generalized null space.

\begin{exercises}
  \recommended \item
    给出零变换与恒等变换的值域空间链和零空间链。
    \begin{answer}
      对于零变换，无论空间是什么，值域空间链都是
      \( V\supset\set{\vec{0}}=\set{\vec{0}}=\cdots\, \)，
      而零空间链是 \( \set{\vec{0}}\subset V=V=\cdots\, \)。
      对于恒等变换，两条链分别是
      \( V=V=V=\cdots \) 与
      \( \set{\vec{0}}=\set{\vec{0}}=\cdots\, \)。
    \end{answer}
  \recommended\item
     对每个映射，
     给出值域空间链与零空间链，
     以及广义值域空间与
     广义零空间。
     \begin{exparts}
       \partsitem $\map{t_0}{\polyspace_2}{\polyspace_2}$,
         $a+bx+cx^2\mapsto b+cx^2$
       \partsitem $\map{t_1}{\Re^2}{\Re^2}$,
         \begin{equation*}
           \colvec{a \\ b}\mapsto\colvec{0 \\ a}
         \end{equation*}
       \partsitem $\map{t_2}{\polyspace_2}{\polyspace_2}$,
         $a+bx+cx^2\mapsto b+cx+ax^2$
       \partsitem $\map{t_3}{\Re^3}{\Re^3}$,
         \begin{equation*}
           \colvec{a \\ b \\ c}\mapsto\colvec{a \\ a \\ b}
         \end{equation*}
     \end{exparts}
     \begin{answer}
       \begin{exparts}
         \partsitem 将 $t_0$ 迭代两次
           $a+bx+cx^2\mapsto b+cx^2\mapsto cx^2$
           得到
           \begin{equation*}
             a+bx+cx^2\mapsunder{t_0^2}cx^2
           \end{equation*}
           而任何更高的迭代都是同一个映射。
           于是，虽然 $\rangespace{t_0}$ 是没有一次项的
           二次多项式构成的空间
           $\set{p+rx^2\suchthat p,r\in \C}$，
           且
           $\rangespace{t_0^2}$ 是纯二次多项式构成的空间
           $\set{rx^2\suchthat r\in \C}$，
           但链就在这里稳定下来：
           $\genrangespace{t_0}=\set{rx^2\suchthat r\in \C}$。

           至于零空间，
           $\nullspace{t_0}$ 是仅含一次项的
           多项式构成的空间 $\set{qx\suchthat q\in \C}$，
           而
           $\nullspace{t_0^2}$ 是一次多项式构成的空间
           $\set{p+qx\suchthat p,q\in \C}$。
           到此为止：$\gennullspace{t_0}=\nullspace{t_0^2}$。
         \partsitem 二次幂
           \begin{equation*}
             \colvec{a \\ b}
              \mapsunder{t_1}\colvec{0 \\ a}
              \mapsunder{t_1}\colvec{0 \\ 0}
           \end{equation*}
           是零映射。
           因此，值域空间链
           \begin{equation*}
             \Re^2
               \supset\set{\colvec{0 \\ p}\suchthat p\in\C}
               \supset\set{\zero}
               =\cdots
           \end{equation*}
           与零空间链
           \begin{equation*}
             \set{\zero}
               \subset\set{\colvec{q \\ 0}\suchthat q\in\C}
               \subset\Re^2
               =\cdots
           \end{equation*}
           的长度都是二。
           广义值域空间是平凡子空间，而
           广义零空间是整个空间。
         \partsitem 这个映射的迭代循环往复
           \begin{equation*}
             a+bx+cx^2
               \mapsunder{t_2} b+cx+ax^2
               \mapsunder{t_2} c+ax+bx^2
               \mapsunder{t_2} a+bx+cx^2
               \;\cdots
           \end{equation*}
           并且值域空间链与零空间链都是平凡的。
           \begin{equation*}
             \polyspace_2=\polyspace_2=\cdots
              \qquad
              \set{\zero}=\set{\zero}=\cdots
           \end{equation*}
           于是，显然，
           广义空间是 $\genrangespace{t_2}=\polyspace_2$
           与 $\gennullspace{t_2}=\set{\zero}$。
         \partsitem 我们有
           \begin{equation*}
             \colvec{a \\ b \\ c}
                \mapsto\colvec{a \\ a \\ b}
                \mapsto\colvec{a \\ a \\ a}
                \mapsto\colvec{a \\ a \\ a}
                \mapsto\cdots
           \end{equation*}
           于是值域空间链
           \begin{equation*}
             \Re^3
               \supset\set{\colvec{p \\ p \\ r}\suchthat p,r\in\C}
               \supset\set{\colvec{p \\ p \\ p}\suchthat p\in\C}
               =\cdots
           \end{equation*}
           以及零空间链
           \begin{equation*}
             \set{\zero}
                \subset\set{\colvec{0 \\ 0 \\ r}\suchthat r\in \C}
                \subset\set{\colvec{0 \\ q \\ r}\suchthat q,r\in \C}
                =\cdots
           \end{equation*}
           的长度都是二。
           广义空间就是上面每条链中最后所示的那个空间。
       \end{exparts}
      \end{answer}
  \item
    证明函数的复合满足结合律
    \( \composed{(\composed{t}{t})}{t}=\composed{t}{(\composed{t}{t})} \)，
    因而我们写 $t^3$ 时无须指明结合方式。
    \begin{answer}
      两者都把 \( x\mapsto t(t(t(x))) \)。
    \end{answer}
  \item \label{exer:PropSubspStrictLowerDimen}
    验证子空间的维数必定小于或等于
    其母空间的维数。
    验证若子空间是真子空间（子空间不等于
    母空间），则维数严格更小。
    \textit{（这用于
             \nearbylemma{le:RangeAndNullChains} 的证明中。）}
    \begin{answer}
      回忆一下，若 $W$ 是 $V$ 的子空间，则我们可以把
      $W$ 的任意一组基 $B_W$
      扩充为 $V$ 的一组基 $B_V$。
      由此第一句话立即可得。
      第二句话也不难：~$W$ 是 $B_W$ 的张成空间，
      若 $W$ 是真子空间，则 $V$ 不是 $B_W$ 的张成空间，
      所以 $B_V$ 至少比 $B_W$ 多一个向量。
    \end{answer}
  \recommended\item
    证明：若变换 $t$ 是非奇异的，则广义值域空间 $\genrangespace{t}$ 是
    整个空间，且广义零空间 $\gennullspace{t}$ 是平凡的。
    反过来“仅当”方向也成立吗？
    \begin{answer}
      “若”与“仅当”都成立。
      线性映射是非奇异的
      当且仅当它保持维数，也就是说，其值域的维数
      等于其定义域的维数。
      对于变换 $\map{t}{V}{V}$，这意味着
      该映射是非奇异的当且仅当它是满射：
      $\rangespace{t}=V$（从而 $\rangespace{t^2}=V$，等等）。
    \end{answer}
  \item \label{exer:RangeAndNullChains}
    验证 \nearbylemma{le:RangeAndNullChains} 中零空间的那一半。
    \begin{answer}
      零空间构成链，因为
      若 $\vec{v}\in\nullspace{t^j}$，则 $t^j(\vec{v})=\zero$，
      且 $t^{j+1}(\vec{v})=t(\,t^j(\vec{v})\,)=t(\zero)=\zero$，
      于是 $\vec{v}\in\nullspace{t^{j+1}}$。

      现在，
      零空间的“进一步”性质由它对值域空间成立这一事实
      结合前面的练习即可得出。
      由于 $\rangespace{t^j}$ 的维数加上
      $\nullspace{t^j}$ 的维数等于起始空间~$V$ 的维数~$n$，
      当值域空间的维数停止下降时，
      零空间的维数也不再变化。
      前面的练习表明，从这一点~$k$ 开始，
      链中的包含关系都不是真包含\Dash 零空间
      彼此相等。
    \end{answer}
  \recommended\item
    给出一个三维空间上的变换的例子，使其值域的维数为二。
    它的零空间是什么？
    对你的例子进行迭代，直到值域空间与零空间稳定。
    \begin{answer}
      （许多例子都正确，这里给出一个。）
      一个例子是实三元组上的平移算子
      \( (x,y,z)\mapsto (0,x,y) \)。
      零空间是所有以两个零开头的三元组。
      该映射在三次迭代后稳定。
     \end{answer}
  \item
      证明线性变换的值域空间与零空间未必不相交。
      它们有可能不相交吗？
      \begin{answer}
        求微分算子
        \( \map{d/dx}{\polyspace_1}{\polyspace_1} \) 的值域空间与零空间相同。
        要看一个二者不相交\Dash
        除零向量外\Dash 的例子，
        可考虑恒等映射，或任意非奇异映射。
      \end{answer}
\end{exercises}

















\subsectionoptional{链}
\index{nilpotent|(}
\noindent\textit{本小节需要用到选读
    小节「子空间的组合」中的内容。}

前一小节表明，随着 \( j \) 的增大，
$\rangespace{t^j}$ 的维数下降，
而 $\nullspace{t^j}$ 的维数上升，
并且此秩与零度在二者之间分摊了
$V$ 的维数。
我们还能说得更深入一些吗？
这两者是否分拆出一个基\Dash 即
\( V=\rangespace{t^j}\directsum\nullspace{t^j} \) 是否成立？

对最小的幂 $j=0$ 答案是肯定的，因为
\( V=\rangespace{t^0}\directsum\nullspace{t^0}=V\directsum\set{\zero} \)。
在另一个极端情形，答案也是肯定的。

\begin{lemma} \label{lem:RestONeToOne}
对任意线性映射 \( \map{t}{V}{V} \)，
函数 \( \map{t}{\genrangespace{t}}{\genrangespace{t}} \)
是单射。
\end{lemma}

\begin{proof}
设 $V$ 的维数为 $n$。
因为 \( \rangespace{t^n}=\rangespace{t^{n+1}} \)，
所以映射 \( \map{t}{\genrangespace{t}}{\genrangespace{t}} \)
是保持维数的同态。
因此，由定理~Three.II.\ref{th:OOHomoEquivalence} 可知它是单射。
\end{proof}

\begin{corollary} \label{GenRngNullDirSumToSp}
设 \( \map{t}{V}{V} \) 是线性变换，
则该空间是直和
\( V=\genrangespace{t}\directsum\gennullspace{t} \)。
也就是说，(1)~\( \dim(V)=\dim(\genrangespace{t})+\dim(\gennullspace{t}) \)
与 (2)~\( \genrangespace{t}\intersection\gennullspace{t}=\set{\zero} \)
均成立。
\end{corollary}

\begin{proof}
设 $V$ 的维数为 $n$。
我们来验证第二句话，它与第一句等价。
第~(1) 条成立，因为任何变换都满足
其秩加上其零度
等于空间的维数，
特别地，这对变换 $t^n$ 也成立。

对于第~(2) 条，假设
\( \vec{v}\in\genrangespace{t}\intersection\gennullspace{t} \)
%               =\rangespace{t^n}\intersection\nullspace{t^n} \)
我们来证明 $\vec{v}=\zero$。
因为 \( \vec{v} \) 属于广义零空间，所以 \( t^n(\vec{v})=\zero \)。
另一方面，由引理
\( \map{t}{\genrangespace{t}}{\genrangespace{t}} \)
是单射，而
单射映射的复合仍是单射，所以
\( \map{t^n}{\genrangespace{t}}{\genrangespace{t}} \) 是单射。
单射线性映射只有把 \( \zero \) 映到
\( \zero \)，故 \( t^n(\vec{v})=\zero \) 就意味着
\( \vec{v}=\zero \)。
\end{proof}

\begin{remark}
严格说来，映射 $\map{t}{V}{V}$ 与
子空间上的映射 \( \map{t}{\genrangespace{t}}{\genrangespace{t}} \)
在广义值域空间不等于 $V$ 时是有区别的，因为二者的定义域
不同。
但这个差别不大，因为
后者就是前者
限制 % \index{restriction}\index{map!restriction}\appendrefs{map restrictions}\spacefactor=1000 %
在 $\genrangespace{t}$ 上
所得的映射。
\end{remark}

对于介于 $j=0$ 与~$j=n$ 之间的幂，
空间 $V$ 可能不是
$\rangespace{t^j}$ 与 $\nullspace{t^j}$ 的直和。
下面的例子表明，这两者可能有非平凡的交。

\begin{example}   \label{FirstNilMap}
考察 \( \C^2 \) 上的一个变换，
它按如下方式作用于标准基的元素。
\begin{equation*}
  \colvec[r]{1 \\ 0}
    \mapsunder{n}
    \colvec[r]{0 \\ 1}
  \quad
  \colvec[r]{0 \\ 1}
    \mapsunder{n}
    \colvec[r]{0 \\ 0}
  \qquad
  N=\rep{n}{\stdbasis_2,\stdbasis_2}=\begin{mat}[r]
    0  &0  \\
    1  &0
  \end{mat}
\end{equation*}
这是一个\definend{移位映射}\index{shift}\index{transformation!shift}，
因为它把各个分量向下移位，最下面的分量被完全移
出了向量。
\begin{equation*}
  \colvec{x \\ y}\mapsto \colvec{0 \\ x}
\end{equation*}
在该基上，这个映射的作用给出一条
\definend{链}。\index{string!of basis vectors}
\begin{equation*}
  \begin{strings}{ccccc}
     \colvec{1 \\ 0} &\mapsto &\colvec{0 \\ 1} &\mapsto &\colvec{0 \\ 0}
  \end{strings}
  \qquad\text{即}\qquad
  \begin{strings}{ccccc}
     \vec{e}_1 &\mapsto &\vec{e}_2 &\mapsto &\zero
  \end{strings}
\end{equation*}
这个映射是让向量既属于值域空间
又属于零空间的一种自然方式；链的图示表明
下面这个向量正是这样一个向量。
\begin{equation*}
  \vec{e}_2=\colvec[r]{0 \\ 1}
\end{equation*}
另外可以观察到，虽然 $n$ 不是零映射，
但函数 $n^2=\composed{n}{n}$
是零映射。
\end{example}

\begin{example}  \label{NilIndexFourOnCFour}
设线性函数 \( \map{\hat{n}}{\C^4}{\C^4} \)
在 \( \stdbasis_4 \) 上的作用由下面的链
给出
\begin{equation*}
  \begin{strings}{ccccccccc}
     \vec{e}_1 &\mapsto &\vec{e}_2
          &\mapsto &\vec{e}_3
          &\mapsto &\vec{e}_4
          &\mapsto &\zero
  \end{strings}
\end{equation*}
则有
\( \rangespace{\hat{n}}\intersection\nullspace{\hat{n}} \) 等于
张成空间 \( \spanof{\set{\vec{e}_4}} \)，
有 \( \rangespace{\hat{n}^2}\intersection\nullspace{\hat{n}^2}=
  \spanof{\set{\vec{e}_3,\vec{e}_4}} \)，
且有 \( \rangespace{\hat{n}^3}\intersection\nullspace{\hat{n}^3}=
    \spanof{\set{\vec{e}_4}} \)。
其矩阵表示除次对角线上的若干个
$1$ 之外全为零。
\begin{equation*}
  \hat{N}=\rep{\hat{n}}{\stdbasis_4,\stdbasis_4}
  =\begin{mat}[r]
    0  &0  &0  &0 \\
    1  &0  &0  &0 \\
    0  &1  &0  &0 \\
    0  &0  &1  &0 \
  \end{mat}
\end{equation*}
虽然 $\hat{n}$ 不是零映射，
$\hat{n}^2$ 与~$\hat{n}^3$ 也都不是，但函数
$\hat{n}^4$ 是零函数。
\end{example}

\begin{example} \label{ThirdNilMap}
变换可以经由不止一条链来作用。
变换 \( t \) 按如下方式作用于基
\( B=\sequence{\vec{\beta}_1,\dots,\vec{\beta}_5} \)：
\begin{equation*}
   \begin{strings}{ccccccc}
    \vec{\beta}_1 &\mapsto &\vec{\beta}_2 &\mapsto &\vec{\beta}_3
        &\mapsto &\zero \\
    \vec{\beta}_4 &\mapsto &\vec{\beta}_5 &\mapsto &\zero
  \end{strings}
\end{equation*}
则比如说，$\vec{\beta}_3$ 位于
其值域空间与零空间的交之中。
从这些链可以清楚看出 $t^3$ 是零映射。
这个映射的表示矩阵除了一些由次对角线上的
$1$ 构成的块以外全为零
\begin{equation*}
  \rep{t}{B,B}=
  \begin{pmat}{rrr|rr}
     0  &0  &0  &0  &0  \\
     1  &0  &0  &0  &0  \\
     0  &1  &0  &0  &0  \\ \hline
     0  &0  &0  &0  &0  \\
     0  &0  &0  &1  &0
  \end{pmat}
\end{equation*}
（其中的线条只是为了在视觉上组织这些块）。
\end{example}

在那些例子中，所有向量最终都被变换成
零。

\begin{definition} \label{def:nilpotent} \index{nilpotent!definition}
\definend{幂零的}变换\index{transformation!nilpotent}%
\index{nilpotent!transformation}
是指某个幂为零映射的变换。
\definend{幂零矩阵}\index{matrix!nilpotent}%
\index{nilpotent!matrix}
是指某个幂为零矩阵的矩阵。
无论是哪种情形，最小的这样的幂称为\definend{幂零指数}。%
\index{nilpotency, index of}\index{index of nilpotency}
\end{definition}

\begin{example}
在\nearbyexample{FirstNilMap}中幂零指数是二。
在\nearbyexample{NilIndexFourOnCFour}中它是四。
在\nearbyexample{ThirdNilMap}中它是三。
\end{example}

\begin{example}
求导映射 \( \map{d/dx}{\polyspace_2}{\polyspace_2} \)
是幂零指数为三的幂零映射，因为任何二次多项式的三阶导数为零。
这个映射的作用由链
$x^2\mapsto 2x\mapsto 2\mapsto 0$
描述，而取基 \( B=\sequence{x^2,2x,2} \)
就给出这个表示。
\begin{equation*}
  \rep{d/dx}{B,B}=
  \begin{mat}[r]
     0  &0  &0  \\
     1  &0  &0  \\
     0  &1  &0
  \end{mat}
\end{equation*}
\end{example}

并非所有幂零矩阵都是除了由次对角线上的
$1$ 构成的块以外全为零的矩阵。

\begin{example} \label{ex:NilMatNotCanon}
利用\nearbyexample{NilIndexFourOnCFour}中的矩阵 $\hat{N}$，
以及这个由四个向量构成的基
\begin{equation*}
  D=\sequence{\colvec[r]{1 \\ 0 \\ 1 \\ 0},
              \colvec[r]{0 \\ 2 \\ 1 \\ 0},
              \colvec[r]{1 \\ 1 \\ 1 \\ 0},
              \colvec[r]{0 \\ 0 \\ 0 \\ 1}}
\end{equation*}
一次换基操作
得到了这个关于 \( D,D \) 的表示。
\begin{equation*}
  \begin{mat}[r]
    1  &0  &1 &0 \\
    0  &2  &1 &0 \\
    1  &1  &1 &0 \\
    0  &0  &0 &1
  \end{mat}
  \begin{mat}[r]
    0  &0  &0 &0 \\
    1  &0  &0 &0 \\
    0  &1  &0 &0 \\
    0  &0  &1 &0
  \end{mat}
  \begin{mat}[r]
    1  &0  &1 &0 \\
    0  &2  &1 &0 \\
    1  &1  &1 &0 \\
    0  &0  &0 &1
  \end{mat}^{-1}\!\!
  =
  \begin{mat}[r]
   -1  &0  &1   &0 \\
   -3  &-2 &5   &0 \\
   -2  &-1  &3  &0 \\
    2  &1   &-2 &0
  \end{mat}
\end{equation*}
新矩阵是幂零的；它的四次幂
是零矩阵。
我们可以用繁琐的计算来验证这一点，也可以
只是注意到
它是幂零的，因为它的四次幂
相似于 \( \hat{N}^4 \)，即零矩阵，
而相似于零矩阵的矩阵只有它自己。
\begin{equation*}
   (P\hat{N}P^{-1})^4
   =P\hat{N}P^{-1}\cdot P\hat{N}P^{-1}\cdot P\hat{N}P^{-1}\cdot P\hat{N}P^{-1}
   =P\hat{N}^4P^{-1}
\end{equation*}
\end{example}

本小节的目标是证明
前面的例题是典型性的，
即每个幂零矩阵都相似于一个除了由次对角线上的
$1$ 构成的块以外全为零的矩阵。

\begin{definition}
设 \( t \) 是 \( V \) 上的幂零变换。
由 \( \vec{v}\in V \) 生成的
\definend{\( t \)-链}\index{string}
是序列
\( \sequence{\vec{v},t(\vec{v}),\ldots,t^{k-1}(\vec{v})} \)，
其中 $t^k(\vec{v})=\zero$。
\definend{\( t \)-链基}\index{basis!string}\index{string!basis}
是由若干条 \( t \)-链拼接而成的基。
\end{definition}

\noindent （这些链只有在互不相交时才能拼接成基，
因为基中不能有重复的向量。）

\begin{example}
这个线性映射 \( \map{t}{\C^3}{\C^3} \)
\begin{equation*}
  \colvec{x \\ y \\ z}\mapsunder{t}\colvec{y \\ z \\ 0}
\end{equation*}
是幂零的，幂零指数为~$3$。
\begin{equation*}
  \colvec{x \\ y \\ z}
  \mapsunder{t}\colvec{y \\ z \\ 0}
  \mapsunder{t}\colvec{z \\ 0 \\ 0}
  \mapsunder{t}\colvec{0 \\ 0 \\ 0}
\end{equation*}
这是一条 $t$-链。
\begin{equation*}
  \sequence{\colvec{0 \\ 0 \\ 1},
            \colvec{0 \\ 1 \\ 0},
            \colvec{1 \\ 0 \\ 0}}
\end{equation*}
因为这个序列是一个基，
所以它是空间~$\C^3$ 的一组 $t$-链基。
\end{example}

当前面的例子中的序列为
$\sequence{\vec{\beta}_1,
             \vec{\beta}_2,
             \vec{\beta}_3}$ 时，
另一条 $t$-链是
$\sequence{\vec{\beta}_2,
             \vec{\beta}_3}$。
但当然，第二个序列并不是 $\C^3$ 的基。
在这个意义上，$t$-链基中的 $t$-链必须是极大的。

\begin{example}
求导线性映射 \( \map{d/dx}{\polyspace_2}{\polyspace_2} \)
是幂零的。
序列
\( \sequence{x^2, 2x, 2} \)
是一条长度为~\( 3 \) 的 \( d/dx \)-链；
特别地，这条链满足 $d/dx(2)=0$ 的要求。
由于它是一个基，该序列
是 \( \polyspace_2\) 的一个 \( d/dx \)-链基。
\end{example}

\begin{example}
在\nearbyexample{ThirdNilMap}中，我们可以把 $t$-链
$\sequence{\vec{\beta}_1,\vec{\beta}_2,\vec{\beta}_3}$ 与
$\sequence{\vec{\beta}_4,\vec{\beta}_5}$
拼接起来，作为 $t$ 的定义域的一组基。
\end{example}

\begin{lemma}  \label{le:LongestTowerIsIndex}
如果一个空间有 \( t \)-链基，那么 $t$ 的幂零指数
等于该基中最长链的长度。
\end{lemma}

\begin{proof}
设该空间有一组 $t$-链基，且 $t$ 的幂零
指数为~$k$。
那么 \( t^k \) 把任何向量都送到 \( \zero \)。
这必须包括
每条链起始的向量，
所以链基中每条链的长度至多为~$k$。

反之，假设该空间有一组 $t$-链基~$B$，其中
所有链的长度都短于~\( k \)。
因为 $t$ 的幂零指数为~$k$，所以存在 \( \vec{v} \)
使得 \( t^{k-1}(\vec{v})\neq\zero \)。
把 $\vec{v}$ 表示为~$B$ 中元素的线性组合，
并施以 \( t^{k-1} \)。
我们假设了 \( t^{k-1} \) 把~$B$ 的每个元素都映到 \( \zero \)。
于是它把线性组合中的每一项都映到 $\zero$，
这与它不把 \( \vec{v} \) 映到 \( \zero \) 矛盾。
\end{proof}

我们将证明每个
幂零映射都有一组相伴的链基，即由互不相交的链组成的基。
% Then our goal theorem, that every
% nilpotent matrix is similar to one that is
% all zeros except for blocks of subdiagonal ones, is immediate,
 %as in \nearbyexample{ThirdNilMap}.

为了看清这一论证的主要思想，设想
我们想构造一个反例：一个幂零的映射，
但没有相伴的由互不相交的链组成的基。
我们可能会想到构造类似
\( \map{t}{\C^5}{\C^5} \) 且作用如下所示的映射。
\begin{center}
  \begin{minipage}{1in}
  \setlength{\unitlength}{1pt}
  $\begin{strings}{ccccc}
     \begin{picture}(4,15)
       \put(0,14){$\vec{e}_1$}
       \put(0,-12){$\vec{e}_2$}
     \end{picture}
     &\begin{picture}(8,10)
       \put(0,8){\rotatebox{-30}{$\mapsto$}}%
       \put(0,-8){\rotatebox{30}{$\mapsto$}}
     \end{picture}
     &\vec{e}_3 &\mapsto &\zero  \\[4ex]
     \vec{e}_4 &\mapsto &\vec{e}_5 &\mapsto &\zero
  \end{strings}$
  \end{minipage}
  \hspace*{3em}
  $\rep{t}{\stdbasis_5,\stdbasis_5}=
  \begin{mat}[r]
    0  &0  &0  &0  &0  \\
    0  &0  &0  &0  &0  \\
    1  &1  &0  &0  &0  \\
    0  &0  &0  &0  &0  \\
    0  &0  &0  &1  &0
  \end{mat}$
\end{center}
但是，所示基不是互不相交的这一事实并不意味着
不存在另一个由互不相交的链组成的基。

为了给这个映射构造这样一个基，
我们首先要找出它的链的数目与长度。
注意到 $t$ 的幂零指数是二。
\nearbylemma{le:LongestTowerIsIndex} 表明，在一组互不相交的链基中
至少有一条链的长度为二。
基元素共有五个，所以如果存在互不相交的链基，
那么这个映射的作用方式必是下面两者之一。
\begin{equation*}
  \begin{strings}{ccccc}
    \vec{\beta}_1 &\mapsto &\vec{\beta}_2 &\mapsto &\zero  \\
    \vec{\beta}_3 &\mapsto &\vec{\beta}_4 &\mapsto &\zero  \\
    \vec{\beta}_5 &\mapsto &\zero
  \end{strings}
  \hspace*{3em}
  \begin{strings}{ccccc}
    \vec{\beta}_1 &\mapsto &\vec{\beta}_2 &\mapsto &\zero  \\
    \vec{\beta}_3 &\mapsto &\zero   \\
    \vec{\beta}_4 &\mapsto &\zero   \\
    \vec{\beta}_5 &\mapsto &\zero
  \end{strings}
\end{equation*}
现在是关键之处。
具有左边那种作用的变换，
其零空间的维数是三，因为被映到零的基向量
正好是这么多。
具有右边那种作用的变换的零空间
维数是四。
利用上面的矩阵表示，我们可以确定两种可能的形状中
哪一个是正确的。
\begin{equation*}
  \nullspace{t}=
  \set{\colvec{x \\ -x \\ z \\ 0 \\ r}\suchthat x,z,r\in\C }
\end{equation*}
它是三维的，
这意味着在上面两种互不相交链基的形式中，\( t \) 的
基是左边那一种。

为了构造~$t$ 的一组链基，首先
从 \( \rangespace{t}\intersection\nullspace{t} \) 中
取 \( \vec{\beta}_2 \) 与 \( \vec{\beta}_4 \)。
\begin{equation*}
  \vec{\beta}_2=\colvec[r]{0 \\ 0 \\ 1 \\ 0 \\ 0}\qquad
  \vec{\beta}_4=\colvec[r]{0 \\ 0 \\ 0 \\ 0 \\ 1}
\end{equation*}
（其他选择也是可能的，只需保证集合
\( \set{\vec{\beta}_2,\vec{\beta}_4} \) 线性无关。）
对于 \( \vec{\beta}_5 \)，从 \( \nullspace{t} \) 中取一个
不在 \( \set{ \vec{\beta}_2,\vec{\beta}_4 } \) 的张成之中的向量。
\begin{equation*}
  \vec{\beta}_5=\colvec[r]{1 \\ -1 \\ 0 \\ 0 \\ 0}
\end{equation*}
最后，取 \( \vec{\beta}_1 \) 与 \( \vec{\beta}_3 \) 使得
\( t(\vec{\beta}_1)=\vec{\beta}_2 \) 且
\( t(\vec{\beta}_3)=\vec{\beta}_4 \)。
\begin{equation*}
  \vec{\beta}_1=\colvec[r]{0 \\ 1 \\ 0 \\ 0 \\ 0}\qquad
  \vec{\beta}_3=\colvec[r]{0 \\ 0 \\ 0 \\ 1 \\ 0}
\end{equation*}
这样，我们就有了一组链基
\( B=\sequence{\vec{\beta}_1,\ldots,\vec{\beta}_5} \)，
而关于这组基，
$t$ 的矩阵具有由次对角线~$1$ 构成的块。
\begin{equation*}
  \rep{t}{B,B}=
  \begin{pmat}{rr|rr|r}
    0  &0  &0  &0  &0  \\
    1  &0  &0  &0  &0  \\  \hline
    0  &0  &0  &0  &0  \\
    0  &0  &1  &0  &0  \\  \hline
    0  &0  &0  &0  &0
  \end{pmat}
\end{equation*}

\begin{theorem}
\label{th:NilMapHasStrBas}
任何幂零变换 $t$ 都有一组与之相伴的 \( t \)-链基。
虽然这组基并不唯一，但链的数目
与长度由 \( t \) 决定。
\end{theorem}

下面的图示说明了这个证明，其中描述了三类
基向量。
若它们属于零空间则画在方框中，
否则画在圆圈中。
\begin{equation*}
   \begin{strings}{ccccccccccccccccccc}
     \digitincirc{3}
         &\mapsto &\digitincirc{1} &\mapsto
         &\cdots &  &  &  &  &  &  &  &\cdots
         &\mapsto &\digitincirc{1}
         &\mapsto &\digitinsq{1} &\mapsto &\zero  \\[.75ex]
     \digitincirc{3}
         &\mapsto &\digitincirc{1} &\mapsto
         &\cdots &  &  &  &  &  &  &  &\cdots
         &\mapsto &\digitincirc{1}
         &\mapsto &\digitinsq{1} &\mapsto &\zero  \\[.75ex]
         &\smash{\vdotswithin{\mapsto}}  \\
     \digitincirc{3}
         &\mapsto &\digitincirc{1} &\mapsto
         &\cdots &
         &\mapsto &\digitincirc{1}
         &\mapsto &\digitinsq{1} &\mapsto &\zero  \\[1ex]
     \digitinsq{2} &\mapsto &\zero \\[.75ex]
         &\smash{\vdotswithin{\mapsto}}  \\
     \digitinsq{2} &\mapsto &\zero
   \end{strings}
\end{equation*}

\begin{proof}
固定一个向量空间 $V$。
我们将对幂零指数作归纳论证。
如果映射 $\map{t}{V}{V}$
的幂零指数为~\( 1 \)，那么它是零映射，且任何基
都是链基：$\vec{\beta}_1\mapsto\zero$，\ldots，
$\vec{\beta}_n\mapsto\zero$。

对于归纳步骤，假设该定理对任何幂零指数从 $1$ 直到且包括~\( k-1 \)（其中 $k>1$）
的变换 $\map{t}{V}{V}$
都成立，
并考虑指数为~$k$ 的情形。

使归纳得以进行的是这样一个观察：
\( t \) 限制在值域空间~\( \rangespace{t} \) 上
也是幂零的，幂零指数为 \( k-1 \)。
于是应用归纳假设，得到
\( \rangespace{t} \) 的一组链基，
其中链的数目与长度
由 \( t \) 决定。
\begin{equation*}
  B=\cat{\cat{\sequence{\vec{\beta}_1,t(\vec{\beta}_1),\dots,
     t^{h_1}(\vec{\beta}_1)}}{
  \cat{\sequence{\vec{\beta}_2,\ldots,t^{h_2}(\vec{\beta}_2)}}}{\cdots}}{
  \sequence{\vec{\beta}_i,\ldots,t^{h_i}(\vec{\beta}_i)} }
\end{equation*}
我们用~$i$ 表示链的条数。
在上面的图示中，这些就是第~\( 1 \) 类向量。

取每条链中的最后一个向量，
就得到交 \( \rangespace{t}\intersection\nullspace{t} \) 的
一组基
\( C=\sequence{t^{h_1}(\vec{\beta}_1),\dots,t^{h_i}(\vec{\beta}_i)} \)。
这是因为 \( \rangespace{t} \) 中的成员映到零当且仅当
它是那些映到
零的基向量的线性组合。
图示中这些向量画成方框里的 \( 1 \)。

现在把 \( C \) 扩充为整个零空间 \( \nullspace{t} \) 的基。
\begin{equation*}
  \hat{C}=\cat{C}{\sequence{\vec{\xi}_1,\dots,\vec{\xi}_p}}
\end{equation*}
虽然 \( \vec{\xi} \) 们并不由 $t$ 唯一
确定，但唯一确定的是它们的
数目：~它等于
\( \nullspace{t} \) 的维数减去
\( \rangespace{t}\intersection\nullspace{t} \) 的维数。
在图示中，$\vec{\xi}$ 们是第~\( 2 \) 类向量，
因而 \( \hat{C} \) 就是画在方框中的向量的集合。

最后，\( \cat{B}{\sequence{\vec{\xi}_1,\dots,\vec{\xi}_p}} \)
是 \( \rangespace{t}+\nullspace{t} \) 的一组基。
这是因为值域空间中的元素与零空间中的元素的
和
可以用 \( B \) 中的元素表示值域空间部分，
而任何来自
\( \nullspace{t}-\rangespace{t} \) 的部分则用 $\xi$ 们表示。
注意到
\begin{align*}
  \dim\big(\rangespace{t}+\nullspace{t}\big)
  &=
  \dim (\rangespace{t})+\dim (\nullspace{t})
    -\dim\big(\rangespace{t}\intersection\nullspace{t}\big)  \\
  &=
  \rank (t)+\nullity (t)-i          \\
  &=
  \dim (V)-i
\end{align*}
于是我们可以把基
\( \cat{B}{\sequence{\vec{\xi}_1,\dots,\vec{\xi}_p}} \)
扩充为整个 \( V \) 的 $t$-链基，
只需再添加~\( i \) 个向量。
为了造出这些向量，
回忆 \( \vec{\beta}_1,\dots,\vec{\beta}_i \) 中的每一个
都属于值域空间 \( \rangespace{t} \)，
因此要用到的向量
\( \vec{v}_1,\dots,\vec{v}_i\in V \) 满足
\( t(\vec{v}_1)=\vec{\beta}_1,\dots,t(\vec{v}_i)=\vec{\beta}_i \)。
验证这一扩充保持线性无关是
\nearbyexercise{exer:NilMapWillHaveStrBas}。
在图示中，\( \vec{v}_1,\dots,\vec{v}_i \) 就是那些 \( 3 \)。
\end{proof}

\begin{corollary} \label{cor:NilpotentMatCanonForm}
\index{nilpotent!canonical form for}
\index{transformation!nilpotent!canonical representative}
\index{canonical form!for nilpotent matrices}
每个幂零矩阵都相似于一个除了由次对角线上的
$1$ 构成的块以外全为零的矩阵。
也就是说，每个幂零映射关于某个基都可以由
这样的矩阵表示。
\end{corollary}

这个形式在如下意义上是唯一的：如果一个幂零矩阵相似于两个
这样的矩阵，那么这两个矩阵只是块的排列顺序不同。
因此，只要我们把块按比如说从最长到最短的顺序排列，
它就是幂零矩阵各相似类的
一个标准形。

\begin{example}
矩阵
\begin{equation*}
  M=\begin{mat}[r]
      1  &-1  \\
      1  &-1
    \end{mat}
\end{equation*}
的幂零指数为二，如下面的计算所示。
\begin{center}
  \begin{tabular}{c|cc}
    \multicolumn{1}{c}{\textit{幂} \( p \)}  &\( M^p \)  &\( \nullspace{M^p}  \)   \\
    \hline
    \( 1 \)
    &\(
       M=\matrixvenlarge{\begin{mat}[r]
         1  &-1  \\
         1  &-1
       \end{mat}}  \)
    &\( \set{\matrixvenlarge{\colvec{x \\ x}}\suchthat
                               x\in\C}  \)   \\
    \( 2 \)
    &\(  M^2=\matrixvenlarge{\begin{mat}[r]
         0  &0   \\
         0  &0
       \end{mat}}  \)
    &\( \C^2  \)
  \end{tabular}
\end{center}
因为该矩阵是 $\nbyn{2}$ 的，它所表示的任何变换
都作用在一个二维空间上。
映射作用一次的零空间
$\nullspace{m}$ 的维数为一，而
作用两次的零空间 $\nullspace{m^2}$ 的维数为二。
于是 $m$ 在某组链基上的作用是
$\vec{\beta}_1\mapsto\vec{\beta}_2\mapsto\zero$，
且该矩阵的标准形如下。
\begin{equation*}
  N=\begin{mat}[r]
    0  &0  \\
    1  &0
  \end{mat}
\end{equation*}

我们可以具体给出这样一组链基，
同时给出见证 $M$ 与~$N$ 矩阵相似的换基矩阵。
假设 $\map{m}{\C^2}{\C^2}$ 满足：\( M \) 关于标准基
表示它。
（我们也可以取 $M$ 为关于某个其他基的表示，
但标准基更方便。）
取 \( \vec{\beta}_2\in\nullspace{m} \)。
再取 \( \vec{\beta}_1 \)
使得 \( m(\vec{\beta}_1)=\vec{\beta}_2 \)。
\begin{equation*}
  \vec{\beta}_2=\colvec[r]{1 \\ 1}
  \qquad
  \vec{\beta}_1=\colvec[r]{1 \\ 0}
\end{equation*}
对于换基矩阵，回忆一下相似示意图。
\begin{equation*}
  \begin{CD}
    \C^2_{\wrt{\stdbasis_2}}      @>m>M>        \C^2_{\wrt{\stdbasis_2}}     \\
    @V\scriptstyle\identity V\scriptstyle PV  @V\scriptstyle\identity V\scriptstyle PV \\
    \C^2_{\wrt{B}}                 @>m>N>         \C^2_{\wrt{B}}
  \end{CD}
\end{equation*}
标准形是 \( \rep{m}{B,B}=PMP^{-1} \)，其中
\begin{equation*}
   P^{-1} %=\bigl(\rep{\identity}{\stdbasis_2,B}\bigr)^{-1}
         =\rep{\identity}{B,\stdbasis_2}
         =\begin{mat}[r]
            1  &1  \\
            0  &1
          \end{mat}
   \qquad
   P=(P^{-1})^{-1}
         =\begin{mat}[r]
            1  &-1  \\
            0  &1
          \end{mat}
\end{equation*}
而矩阵计算的验证是常规的。
\begin{equation*}
  \begin{mat}[r]
    1  &-1  \\
    0  &1
  \end{mat}
  \begin{mat}[r]
    1  &-1  \\
    1  &-1
  \end{mat}
  \begin{mat}[r]
    1  &1  \\
    0  &1
  \end{mat}=
  \begin{mat}[r]
    0  &0  \\
    1  &0
  \end{mat}
\end{equation*}
%which does indeed describe the effect of $m$ on the string basis.
\end{example}

\begin{example} \label{ex:FiveByNilStrBas}
这个矩阵
\begin{equation*}
  \begin{mat}[r]
     0  &0  &0  &0  &0  \\
     1  &0  &0  &0  &0  \\
     -1 &1  &1  &-1 &1  \\
     0  &1  &0  &0  &0  \\
     1  &0  &-1 &1  &-1
  \end{mat}
\end{equation*}
是幂零的，幂零指数为~$3$。
% \newlength{\extramatrixvspace}
% \setlength{\extramatrixvspace}{.75ex}
% \newcommand{\matrixvenlarge}[1]{\raisebox{-0.5\height}{\vbox{
%        \vspace*{\extramatrixvspace}
%        \hbox{$#1$}
%         \vspace*{\extramatrixvspace}
%         }}}
\begin{center}
  \begin{tabular}{c|cc}
    \multicolumn{1}{c}{\textit{幂} \( p \)}  &\( N^p \)  &\( \nullspace{N^p}  \)   \\
    \hline
    \( 1 \)
    &\matrixvenlarge{
         \begin{mat}[r]
         0  &0  &0  &0  &0  \\
         1  &0  &0  &0  &0  \\
         -1 &1  &1  &-1 &1  \\
         0  &1  &0  &0  &0  \\
         1  &0  &-1 &1  &-1
         \end{mat}}
    &\(
        \set{\matrixvenlarge{\colvec{0 \\ 0 \\ u-v \\ u \\ v}}
           \suchthat u,v\in\C}  \)                      \\
    \( 2 \)
    &\(
       \matrixvenlarge{\begin{mat}[r]
         0  &0  &0  &0  &0  \\
         0  &0  &0  &0  &0  \\
         1  &0  &0  &0  &0  \\
         1  &0  &0  &0  &0  \\
         0  &0  &0  &0  &0
       \end{mat}}  \)
    &\(
        \set{\matrixvenlarge{\colvec{0 \\ y \\ z \\ u \\ v}}
                       \suchthat y,z,u,v\in\C}  \)  \\
    \( 3 \)
    &\textit{--零矩阵--}
    &\( \C^5 \)
  \end{tabular}
\end{center}
该表告诉我们任何链基的如下信息：~映射作用一次后的
零空间
维数为~二，所以有两个基向量直接映到零；
作用两次后的零空间维数为四，
所以又有两个基向量在第二次迭代时映到零；
而
作用三次后的零空间维数为~五，所以剩下的
一个基向量经三步映到零。
\begin{equation*}
  \begin{strings}{ccccccc}
    \vec{\beta}_1 &\mapsto &\vec{\beta}_2 &\mapsto &\vec{\beta}_3
       &\mapsto &\zero  \\
    \vec{\beta}_4 &\mapsto &\vec{\beta}_5 &\mapsto &\zero
  \end{strings}
\end{equation*}
为了构造这样一个基，首先从
\( \nullspace{n} \) 中取两个构成线性无关集的向量。
\begin{equation*}
   \vec{\beta}_3=\colvec[r]{0 \\ 0 \\ 1 \\ 1 \\ 0} \quad
   \vec{\beta}_5=\colvec[r]{0 \\ 0 \\ 0 \\ 1 \\ 1}
\end{equation*}
然后添加 \( \vec{\beta}_2,\vec{\beta}_4\in\nullspace{n^2} \)，
使得 \( n(\vec{\beta}_2)=\vec{\beta}_3 \) 且
\( n(\vec{\beta}_4)=\vec{\beta}_5 \)。
\begin{equation*}
   \vec{\beta}_2=\colvec[r]{0 \\ 1 \\ 0 \\ 0 \\ 0} \quad
   \vec{\beta}_4=\colvec[r]{0 \\ 1 \\ 0 \\ 1 \\ 0}
\end{equation*}
最后添加 \( \vec{\beta}_1 \)，
使得 \( n(\vec{\beta}_1)=\vec{\beta}_2 \)。
\begin{equation*}
   \vec{\beta}_1=\colvec[r]{1 \\ 0 \\ 1 \\ 0 \\ 0}
\end{equation*}
\end{example}







\begin{exercises}
   \recommended \item \label{exer:IndNilLftShift}
      \definend{右移}算子的幂零指数是多少？这里它作用在实数三元组构成的空间上：
       \begin{equation*}
          (x,y,z)\mapsto(0,x,y)
       \end{equation*}
       \begin{answer}
         是三。
         说它至少是三，是因为 $\ell^2(\,(1,1,1)\,)=(0,0,1)\neq \zero$。
         说它至多是三，是因为
         $(x,y,z)\mapsto (0,x,y)\mapsto (0,0,x)\mapsto (0,0,0)$。
       \end{answer}
   \recommended \item
     对下面每个链基，给出幂零映射的幂零指数，并给出该幂零映射每次迭代的
     值域空间与零空间的维数。
     \begin{exparts}
       \partsitem $
         \begin{strings}{ccccccc}
            \vec{\beta}_1 &\mapsto &\vec{\beta}_2 &\mapsto &\zero  \\
            \vec{\beta}_3 &\mapsto &\vec{\beta}_4 &\mapsto &\zero
          \end{strings}$
       \partsitem $
         \begin{strings}{ccccccc}
            \vec{\beta}_1 &\mapsto &\vec{\beta}_2 &\mapsto &\vec{\beta}_3
                 &\mapsto &\zero  \\
            \vec{\beta}_4 &\mapsto &\zero \\
            \vec{\beta}_5 &\mapsto &\zero \\
            \vec{\beta}_6 &\mapsto &\zero
          \end{strings}$
       \partsitem $
         \begin{strings}{ccccccccc}
            \vec{\beta}_1 &\mapsto &\vec{\beta}_2 &\mapsto &\vec{\beta}_3
                 &\mapsto &\zero
          \end{strings}$
     \end{exparts}
     并给出该矩阵的标准形。
     \begin{answer}
       \begin{exparts}
         \partsitem 定义域的维数是四。
           该映射的作用是：空间中任意向量
           $c_1\cdot \vec{\beta}_1+c_2\cdot \vec{\beta}_2
             +c_3\cdot \vec{\beta}_3+c_4\cdot \vec{\beta}_4$
           被送到
           $c_1\cdot \vec{\beta}_2+c_2\cdot \zero
             +c_3\cdot \vec{\beta}_4+c_4\cdot \zero
            =c_1\cdot \vec{\beta}_3+c_3\cdot\vec{\beta}_4$。
           映射的第一次作用把两个基向量 $\vec{\beta}_2$ 与
           $\vec{\beta}_4$ 送到零，
           因此零空间的维数为~二，值域空间的维数也为~二。
           第二次作用后，四个基向量全部被送到零，
           于是二次幂的零空间维数为~四，而二次幂的值域空间维数为零。
           因此幂零指数是二。
           这就是标准形。
           \begin{equation*}
             \begin{mat}[r]
               0  &0  &0  &0  \\
               1  &0  &0  &0  \\
               0  &0  &0  &0  \\
               0  &0  &1  &0
             \end{mat}
           \end{equation*}
         \partsitem 该映射定义域的维数是六。
           对一次幂而言，零空间的维数是四，值域空间的维数是二。
           对二次幂而言，零空间的维数是五，值域空间的维数是一。
           第三次迭代得到的零空间维数为六，值域空间维数为零。
           幂零指数是三，这就是标准形。
           \begin{equation*}
             \begin{mat}[r]
               0  &0  &0  &0  &0  &0  \\
               1  &0  &0  &0  &0  &0  \\
               0  &1  &0  &0  &0  &0  \\
               0  &0  &0  &0  &0  &0  \\
               0  &0  &0  &0  &0  &0  \\
               0  &0  &0  &0  &0  &0
             \end{mat}
           \end{equation*}
         \partsitem 定义域的维数是三，幂零指数是三。
           一次幂的零空间维数为一，值域空间维数为二。
           二次幂的零空间维数为二，值域空间维数为一。
           最后，三次幂的零空间维数为三，值域空间维数为零。
           这是标准形矩阵。
           \begin{equation*}
             \begin{mat}[r]
               0  &0  &0  \\
               1  &0  &0  \\
               0  &1  &0
             \end{mat}
          \end{equation*}
      \end{exparts}
    \end{answer}
  \item
    判断下列矩阵中哪些是幂零的。
    \begin{exparts*}
      \partsitem
        $\begin{mat}[r]
           -2  &4  \\
           -1  &2
         \end{mat}$
       \partsitem
         $\begin{mat}[r]
           3  &1  \\
           1  &3
         \end{mat}$
       \partsitem
         $\begin{mat}[r]
           -3  &2  &1  \\
           -3  &2  &1  \\
           -3  &2  &1
         \end{mat}$
       \partsitem
         $\begin{mat}[r]
            1  &1  &4  \\
            3  &0  &-1 \\
            5  &2  &7
         \end{mat}$
       \partsitem
         $\begin{mat}[r]
            45  &-22  &-19  \\
            33  &-16  &-14  \\
            69  &-34  &-29
         \end{mat}$
     \end{exparts*}
     \begin{answer}
       由 \nearbylemma{le:RangeAndNullChains}，零度会在第 $n$ 次迭代时增长到尽可能大，其中 $n$ 是定义域的维数。
       于是，对 $\nbyn{2}$ 矩阵，
       我们只需检查其平方是否为零矩阵。
       对 $\nbyn{3}$ 矩阵，只需检查其立方。
       \begin{exparts}
         \partsitem 是，该矩阵是幂零的，因为
           它的平方是零矩阵。
         \partsitem 否，它的平方不是零矩阵。
           \begin{equation*}
             \begin{mat}[r]
                3  &1  \\
                1  &3
             \end{mat}^2
             =\begin{mat}[r]
                10  &6  \\
                6   &10
             \end{mat}
           \end{equation*}
         \partsitem 是，它的立方是零矩阵。
           事实上，它的平方就是零。
         \partsitem 否，它的三次幂不是零矩阵。
           \begin{equation*}
             \begin{mat}[r]
               1  &1  &4  \\
               3  &0  &-1 \\
               5  &2  &7
             \end{mat}^3
             =\begin{mat}[r]
               206  &86  &304  \\
                26  &8   &26   \\
               438  &180 &634
             \end{mat}
           \end{equation*}
         \partsitem 是，该矩阵的立方是零矩阵。
       \end{exparts}
       要看出第二与第四个矩阵不是幂零的，另一种方法是注意它们是非奇异的。
     \end{answer}
   \recommended \item
     求该矩阵的标准形。
     \begin{equation*}
       \begin{mat}[r]
         0  &1  &1  &0  &1  \\
         0  &0  &1  &1  &1  \\
         0  &0  &0  &0  &0  \\
         0  &0  &0  &0  &0  \\
         0  &0  &0  &0  &0
       \end{mat}
     \end{equation*}
     \begin{answer} 计算表格
       \begin{center}
           \begin{tabular}{c|cc}
              \multicolumn{1}{c}{\( p \)}  &\( N^p \)  &\( \nullspace{N^p}  \) \\
              \hline
              \( 1 \)
                &\matrixvenlarge{\begin{mat}[r]
                    0  &1  &1  &0  &1  \\
                    0  &0  &1  &1  &1  \\
                    0  &0  &0  &0  &0  \\
                    0  &0  &0  &0  &0  \\
                    0  &0  &0  &0  &0
                   \end{mat}}
                &\( \set{\matrixvenlarge{\colvec{r \\ u \\ -u-v \\ u \\ v}}
                               \suchthat r,u,v\in\C}  \) \\
              \( 2 \)
                &\matrixvenlarge{\begin{mat}[r]
                    0  &0  &1  &1  &1  \\
                    0  &0  &0  &0  &0  \\
                    0  &0  &0  &0  &0  \\
                    0  &0  &0  &0  &0  \\
                    0  &0  &0  &0  &0
                   \end{mat}}
                &\( \set{\matrixvenlarge{\colvec{r \\ s \\ -u-v \\ u \\ v}}
                               \suchthat r,s,u,v\in\C}  \) \\
             \( 2 \)
                &\textit{--零矩阵--}
                &\( \C^5 \)
           \end{tabular}
         \end{center}
         由此得到对链基的如下要求：~有三个基向量直接被映射到零，
         另有一个基向量经第二次作用被映射到零，
         最后一个基向量
         则经第三次作用被映射到零。
         因此，链基具有如下形式。
         \begin{equation*}
           \begin{strings}{ccccccc}
             \vec{\beta}_1 &\mapsto &\vec{\beta}_2 &\mapsto
               &\vec{\beta}_3 &\mapsto &\zero               \\
             \vec{\beta}_4 &\mapsto &\zero                  \\
             \vec{\beta}_5 &\mapsto &\zero
           \end{strings}
         \end{equation*}
         据此立得标准形。
         \begin{equation*}
           \begin{mat}[r]
             0  &0  &0  &0  &0  \\
             1  &0  &0  &0  &0  \\
             0  &1  &0  &0  &0  \\
             0  &0  &0  &0  &0  \\
             0  &0  &0  &0  &0
           \end{mat}
         \end{equation*}
     \end{answer}
   \recommended \item
     考虑 \nearbyexample{ex:FiveByNilStrBas} 中的矩阵。
     \begin{exparts}
       \partsitem 利用该映射在链基上的作用给出标准形。
       \partsitem 求出把该矩阵化为标准形的换基矩阵。
       \partsitem 利用前一小题的答案检验第一小题的答案。
     \end{exparts}
     \begin{answer}
       \begin{exparts}
         \partsitem 标准形含一个 $\nbyn{3}$ 块与一个
           $\nbyn{2}$ 块
           \begin{equation*}
             \begin{pmat}{rrr|rr}
               0  &0  &0  &0  &0  \\
               1  &0  &0  &0  &0  \\
               0  &1  &0  &0  &0  \\ \hline
               0  &0  &0  &0  &0  \\
               0  &0  &0  &1  &0  \\
             \end{pmat}
           \end{equation*}
           分别对应于基中长度为三的链与长度为二的
           链。
         \partsitem 设 $N$ 是底层映射关于标准基的表示。
           设 $B$ 是我们即将换到的基。
           由相似图
           \begin{equation*}
             \begin{CD}
               \C^2_{\wrt{\stdbasis_2}}
                   @>n>N>
                   \C^2_{\wrt{\stdbasis_2}}     \\
              @V\scriptstyle\identity V\scriptstyle PV
                  @V\scriptstyle\identity V\scriptstyle PV \\
              C^2_{\wrt{B}}
                  @>n>>
              \C^2_{\wrt{B}}
            \end{CD}
          \end{equation*}
          可知标准形矩阵为 $PNP^{-1}$，其中
          \begin{equation*}
            P^{-1}
            =\rep{\identity}{B,\stdbasis_5}
            =\begin{mat}[r]
               1 &0 &0 &0 &0 \\
               0 &1 &0 &1 &0 \\
               1 &0 &1 &0 &0 \\
               0 &0 &1 &1 &1 \\
               0 &0 &0 &0 &1
             \end{mat}
          \end{equation*}
          而 $P$ 是它的逆。
          \begin{equation*}
            P=\rep{\identity}{\stdbasis_5,B}
             =(P^{-1})^{-1}
            =\begin{mat}[r]
               1 &0 &0 &0 &0 \\
              -1 &1 &1 &-1&1 \\
              -1 &0 &1 &0 &0 \\
               1 &0 &-1&1 &-1\\
               0 &0 &0 &0 &1
             \end{mat}
          \end{equation*}
         \partsitem 检验它的计算是例行的。
       \end{exparts}
     \end{answer}
   \recommended \item
     下列每个矩阵都是幂零的。
     \begin{exparts*}
       \partsitem \(
         \begin{mat}[r]
           1/2  &-1/2  \\
           1/2  &-1/2
         \end{mat}        \)
       \partsitem \(
         \begin{mat}[r]
           0  &0  &0  \\
           0  &-1 &1  \\
           0  &-1 &1
         \end{mat}        \)
       \partsitem \(
         \begin{mat}[r]
          -1  &1  &-1 \\
           1  &0  &1  \\
           1  &-1 &1
         \end{mat}        \)
     \end{exparts*}
     把每个矩阵化为标准形。
     \begin{answer}
       \begin{exparts}
       \partsitem 计算
         \begin{center}
           \begin{tabular}{c|cc}
              \multicolumn{1}{c}{\( p \)}  &\( N^p \)  &\( \nullspace{N^p} \) \\
             \hline
              \( 1 \)
                &\matrixvenlarge{\begin{mat}[r]
                     1/2  &-1/2  \\
                     1/2  &-1/2
                   \end{mat}}
                &\( \set{\matrixvenlarge{\colvec{u \\ u}}
                               \suchthat u\in\C}  \) \\
             \( 2 \)
                &\textit{--零矩阵--}
                &\( \C^2 \)
           \end{tabular}
         \end{center}
         表明由该矩阵表示的任何映射
         必须这样作用在链基上
         \begin{equation*}
           \begin{strings}{ccccccc}
             \vec{\beta}_1 &\mapsto &\vec{\beta}_2  &\mapsto &\zero
           \end{strings}
         \end{equation*}
         因为作用一次后零空间的维数为~一，
         且恰好有一个基向量 $\vec{\beta}_2$ 被映射到零。
         因此，这个关于
         $\sequence{\vec{\beta}_1,\vec{\beta}_2}$ 的表示就是标准形。
         \begin{equation*}
           \begin{mat}[r]
             0    &0   \\
             1    &0
           \end{mat}
         \end{equation*}
       \partsitem 这里的计算与前一题类似。
         \begin{center}
           \begin{tabular}{c|cc}
              \multicolumn{1}{c}{\( p \)}  &\( N^p \)  &\( \nullspace{N^p} \) \\
              \hline
              \( 1 \)
                &\matrixvenlarge{\begin{mat}[r]
                     0  &0  &0  \\
                     0  &-1 &1  \\
                     0  &-1 &1
                   \end{mat}}
                &\( \set{\matrixvenlarge{\colvec{u \\ v \\ v}}
                               \suchthat u,v\in\C}  \) \\
            \( 2 \)
                &\textit{--零矩阵--}
                &\( \C^3 \)
           \end{tabular}
        \end{center}
        表格表明链基形如
         \begin{equation*}
           \begin{strings}{ccccccc}
             \vec{\beta}_1 &\mapsto &\vec{\beta}_2 &\mapsto &\zero  \\
             \vec{\beta}_3 &\mapsto &\zero
           \end{strings}
         \end{equation*}
         因为映射作用一次后零空间的
         维数为二\Dash
         --$\vec{\beta}_2$ 与 $\vec{\beta}_3$ 都被送到零\Dash 而
         再迭代一次又使另一个向量
         被送到零。
       \partsitem 计算
         \begin{center}
           \begin{tabular}{c|cc}
              \multicolumn{1}{c}{\( p \)}  &\( N^p \)  &\( \nullspace{N^p} \) \\
              \hline
              \( 1 \)
                &\matrixvenlarge{\begin{mat}[r]
                     -1  &1  &-1  \\
                      1  &0  &1   \\
                      1  &-1 &1
                   \end{mat}}
                &\( \set{\matrixvenlarge{\colvec{u \\ 0 \\ -u}}
                               \suchthat u\in\C}  \) \\
              \( 2 \)
                &\matrixvenlarge{\begin{mat}[r]
                      1  &0  &1   \\
                      0  &0  &0   \\
                     -1  &0  &-1
                   \end{mat}}
                &\( \set{\matrixvenlarge{\colvec{u \\ v \\ -u}}
                               \suchthat u,v\in\C}  \) \\
             \( 3 \)
                &\textit{--零矩阵--}
                &\( \C^3 \)
           \end{tabular}
         \end{center}
         表明由该基表示的任何映射必须这样作用在
         链基上。
         \begin{equation*}
           \begin{strings}{ccccccc}
             \vec{\beta}_1 &\mapsto &\vec{\beta}_2 &\mapsto
               &\vec{\beta}_3 &\mapsto &\zero
           \end{strings}
         \end{equation*}
         因此，这就是标准形。
         \begin{equation*}
           \begin{mat}[r]
             0    &0   &0   \\
             1    &0   &0   \\
             0    &1   &0
           \end{mat}
         \end{equation*}
       \end{exparts}
      \end{answer}
   \item
     描述用幂零矩阵的标准形矩阵左乘或右乘所产生的效果。
     \begin{answer}
       下面几个例子
       \begin{equation*}
         \begin{mat}[r]
           0  &0  \\
           1  &0
         \end{mat}
         \begin{mat}
           a  &b  \\
           c  &d
         \end{mat}
         =
         \begin{mat}
           0  &0  \\
           a  &b
         \end{mat}
         \qquad
         \begin{mat}[r]
           0  &0  &0 \\
           1  &0  &0 \\
           0  &1  &0
         \end{mat}
         \begin{mat}
           a  &b  &c \\
           d  &e  &f \\
           g  &h  &i
         \end{mat}
         =
         \begin{mat}
           0  &0  &0 \\
           a  &b  &c \\
           d  &e  &f
         \end{mat}
       \end{equation*}
       表明用一块次对角线上的 $1$ 左乘一个矩阵，
       会把该矩阵的行向下移。
       而不同的块
       \begin{equation*}
         \begin{mat}[r]
           0  &0  &0  &0  \\
           1  &0  &0  &0  \\
           0  &0  &0  &0  \\
           0  &0  &1  &0
         \end{mat}
         \begin{mat}
           a  &b  &c  &d  \\
           e  &f  &g  &h  \\
           i  &j  &k  &l  \\
           m  &n  &o  &p
         \end{mat}
         =
         \begin{mat}
           0  &0  &0  &0  \\
           a  &b  &c  &d  \\
           0  &0  &0  &0  \\
           i  &j  &k  &l
         \end{mat}
       \end{equation*}
       则会把矩阵的不同部分向下移。

       右乘对列做的是类似的事情。
       参见 \nearbyexercise{exer:IndNilLftShift}。
     \end{answer}
   \item
     幂零性在相似下不变吗？
     也就是说，与幂零矩阵相似的矩阵必定也是幂零的吗？
     若是，指数也相同吗？
     \begin{answer}
       是。
       把 \nearbyexample{ex:NilMatNotCanon} 的最后一句推广即可。
       至于指数，同样的最后一句表明新矩阵的指数小于或等于 $\hat{N}$ 的指数，而把
       两个矩阵的角色互换便得到另一方向的不等式。

       这个问题的另一种回答方式是证明：一个矩阵是幂零的当且仅当与之相关的映射是幂零的，且
       指数相同。
       于是，由于相似矩阵表示同一个映射，结论
       随之而来。
       这就是下面的 \nearbyexercise{exer:MatNilIffMapNil}。
     \end{answer}
   \recommended \item
     证明幂零矩阵唯一的特征值是零。
     \begin{answer}
       注意标准形的幂零矩阵只有
       零特征值；例如，这个下三角矩阵
       \begin{equation*}
          \begin{mat}
            -x  &0  &0  \\
             1  &-x &0  \\
             0  &1  &-x
          \end{mat}
       \end{equation*}
       的行列式是 \( (-x)^3 \)，它的唯一的根是零。
       而相似矩阵有相同的特征值，且每个幂零
       矩阵都与一个标准形的矩阵相似。

       看出这一点的另一种方式是注意：幂零矩阵在若干次迭代后把所有
       向量都送到零，但这与在特征子空间上的作用 $\vec{v}\mapsto \lambda\vec{v}$
       相矛盾，除非
       $\lambda$ 为零。
     \end{answer}
   \item
     在二维空间上存在指数为三的幂零变换吗？
     \begin{answer}
       不存在，由 \nearbylemma{le:RangeAndNullChains}，对二维
       空间上的映射，零度
       在第二次迭代时就会增长到尽可能大。
     \end{answer}
   \item
     在 \nearbytheorem{th:NilMapHasStrBas} 的证明中，为什么归纳基础的情形不是幂零指数为零？
     \begin{answer}
       一个变换的幂零指数为零，仅当开启链的
       向量
       必须是 $\zero$ 时，也就是说，仅当 $V$ 是平凡空间时。
     \end{answer}
   \recommended \item
     设 \( \map{t}{V}{V} \) 是线性变换，并设
     \( \vec{v}\in V \) 满足 \( t^k(\vec{v})=\zero \) 但
     \( t^{k-1}(\vec{v})\neq\zero \)。
     考虑 $t$-链
     $\sequence{\vec{v},t(\vec{v}),\dots,t^{k-1}(\vec{v})}$。
     \begin{exparts}
       \partsitem 证明 \( t \) 是该链中向量组张成的空间上的变换，
         也就是说，证明 \( t\) 限制在该张成空间上的值域
         是该张成空间的子集。
         我们说这个张成空间是 \definend{\( t \)-不变的}
         子空间。\index{invariant subspace}
       \partsitem 证明该限制是幂零的。
       \partsitem 证明该 $t$-链
         线性无关，从而是其张成空间的基。
       \partsitem 用该 $t$-链基表示这个限制映射。
     \end{exparts}
     \begin{answer}
       \begin{exparts}
         \partsitem 张成空间中的任意成员 $\vec{w}$ 都是
           一个线性组合
           $\vec{w}=c_0\cdot \vec{v}+c_1\cdot t(\vec{v})+\dots
            +c_{k-1}\cdot t^{k-1}(\vec{v})$。
           于是，由映射的线性性，
           $t(\vec{w})=c_0\cdot t(\vec{v})+c_1\cdot t^2(\vec{v})+\dots
            +c_{k-2}\cdot t^{k-1}(\vec{v})+c_{k-1}\cdot \zero$
           也在这个张成空间中。
         \partsitem 前一小题中的操作迭代 $k$ 次后，
           会得到零的线性组合。
         \partsitem 若 \( \vec{v}=\zero \)，则该集合是空集，按定义
           线性无关。
           否则，写 \( c_1\vec{v}+\dots+c_{k-1}t^{k-1}(\vec{v})=\zero \)
           并对两边施加 \( t^{k-1} \)。
           右边得到 \( \zero \)，而左边得到
           \( c_1t^{k-1}(\vec{v}) \)；于是 \( c_1=0 \)。
           如此继续，对两边施加 \( t^{k-2} \)，
           依此类推。
         \partsitem 当然，$t$ 把这条长为 $k$ 的单一 $t$-链作为基作用在其张成空间上。
           \begin{equation*}
             \begin{mat}
               0 &0 &0 &0 &\ldots &0 &0 \\
               1 &0 &0 &0 &\ldots &0 &0 \\
               0 &1 &0 &0 &\ldots &0 &0 \\
               0 &0 &1 &0 &       &0 &0 \\
                 &  &  &\ddots          \\
               0 &0 &0 &0 &       &1 &0 \\
             \end{mat}
           \end{equation*}
       \end{exparts}
     \end{answer}
   \item \label{exer:NilMapWillHaveStrBas}
     完成 \nearbytheorem{th:NilMapHasStrBas} 的证明。
     \begin{answer}
       我们必须证明
       \( B\union\hat{C}\union\set{\vec{v}_1,\dots,\vec{v}_j} \) 线性
       无关，其中 \( B \) 是
       \( \rangespace{t} \) 的 \( t \)-链基，\( \hat{C} \) 是
       \( \nullspace{t} \) 的基，且
       \( t(\vec{v}_1)=\vec{\beta}_1,\dots,t(\vec{v}_i)=\vec{\beta}_i \)。
       写出
       \begin{equation*}
         \zero=c_{1,-1}\vec{v}_1+c_{1,0}\vec{\beta}_1
                +c_{1,1}t(\vec{\beta}_1)+\dots+
                c_{1,{h_1}}t^{h_1}(\vec{\vec{\beta}}_1)
                +c_{2,-1}\vec{v}_2+\dots+c_{j,h_i}t^{h_i}(\vec{\beta_i})
       \end{equation*}
       并施加 \( t \)。
       \begin{multline*}
         \zero=c_{1,-1}\vec{\beta}_1+c_{1,0}t(\vec{\beta}_1)
                +\dots+
                c_{1,h_1-1}t^{h_1}(\vec{\vec{\beta}}_1)+c_{1,h_1}\zero      \\
             +c_{2,-1}\vec{\beta}_2+\cdots+c_{i,h_i-1}t^{h_i}(\vec{\beta_i})
             +c_{i,h_i}\zero
       \end{multline*}
       由于 \( B\union\hat{C} \)
       是一组基，可断定系数 \( c_{1,-1},\dots,c_{1,h_i-1},
       c_{2,-1},\dots,c_{i,h_i-1} \) 全为零。
       代回第一个行间公式，可知
       其余系数也为零。
     \end{answer}
   \item \label{exer:MatNilIffMapNil}
     证明 \nearbydefinition{def:nilpotent} 中给出的「幂零变换」与「幂零矩阵」这两个说法彼此吻合：~
     一个映射是幂零的当且仅当它由一个
     幂零矩阵表示。
     （是说一个变换是幂零的当且仅当存在一组基，使得该映射关于该基的表示是
     幂零矩阵，还是说任何表示都是幂零矩阵？）
     \begin{answer}
       对任意基 $B$，
       变换~$n$ 是幂零的当且仅当
       $N=\rep{n}{B,B}$ 是幂零矩阵。
       这是因为只有零矩阵才表示零映射，
       于是 \( n^j \) 是零映射当且仅当 \( N^j \)
       是零矩阵。
     \end{answer}
   \item
     设 \( T \) 是指数为四的幂零变换。
     \( T^3 \) 的值域空间最大能有多大？
     \begin{answer}
       它可以是任何大于或等于一的维数。
       要构造一个指数为四的幂零变换，
       其立方有维数为~$k$ 的值域空间，取一个向量空间、
       该空间的一组基，以及一个按如下方式
       作用于这组基的变换。
       \begin{equation*}
          \begin{strings}{ccccccccc}
             \vec{\beta}_1 &\mapsto &\vec{\beta}_2 &\mapsto &\vec{\beta}_3
               &\mapsto &\vec{\beta}_4 &\mapsto &\zero  \\
             \vec{\beta}_5 &\mapsto &\vec{\beta}_6 &\mapsto &\vec{\beta}_7
               &\mapsto &\vec{\beta}_8 &\mapsto &\zero  \\
             &\vdotswithin{\vec{\beta}_{4k-3}}                       \\
             \vec{\beta}_{4k-3} &\mapsto &\vec{\beta}_{4k-2}
               &\mapsto &\vec{\beta}_{4k-1}
               &\mapsto &\vec{\beta}_{4k} &\mapsto &\zero  \\
             &\vdotswithin{\vec{\beta}_{4k-3}}              \\
             \multicolumn{8}{l}{%
               \text{\textit{--可能还有别的更短的链--}}}
          \end{strings}
       \end{equation*}
       所以 $T^3$ 的值域空间的维数可以任意大。
       它最小为一\Dash 这里
       必须至少有一条链，否则该映射的幂零指数
       就不会是四。
     \end{answer}
   \item
     回忆相似矩阵有相同的特征值。
     证明其逆命题不成立。
     \begin{answer}
       下面两个矩阵的特征值都只有零
       \begin{equation*}
         \begin{mat}
           0  &0  \\
           0  &0
         \end{mat}
         \qquad
         \begin{mat}
           0  &0  \\
           1  &0
         \end{mat}
       \end{equation*}
       但它们并不相似（它们有不同的标准形
       代表元，即它们自身）。
     \end{answer}
   \item \nearbylemma{lem:RestONeToOne} 表明，对任意线性
     变换 \( \map{t}{V}{V} \)，
     限制 \( \map{t}{\genrangespace{t}}{\genrangespace{t}} \)
     都是单射。
     证明它也是满射，从而是一个自同构。
     它必须是恒等映射吗？
     \begin{answer}
       由 \nearbylemma{le:RangeAndNullChains}，它是满射。
       它不必是恒等映射：考虑这个映射 $\map{t}{\Re^2}{\Re^2}$。
       \begin{equation*}
         \colvec{x \\ y}\mapsunder{t}\colvec{y \\ x}
       \end{equation*}
       对该映射有 $\genrangespace{t}=\Re^2$，而 $t$ 不是恒等映射。
     \end{answer}
   \item
     证明幂零矩阵相似于一个除上对角线上的 $1$ 组成的块以外全为零的矩阵。
     \begin{answer}
       对链基作一个简单的重排即可。
       例如，与这个链基相关的映射
       \begin{equation*}
          \begin{strings}{ccccccc}
             \vec{\beta}_1 &\mapsto &\vec{\beta}_2 &\mapsto &\zero
          \end{strings}
       \end{equation*}
       关于
       $B=\sequence{\vec{\beta}_1,\vec{\beta}_2}$ 由这个矩阵表示
       \begin{equation*}
         \begin{mat}
           0  &0  \\
           1  &0
         \end{mat}
       \end{equation*}
       但关于
       $B=\sequence{\vec{\beta}_2,\vec{\beta}_1}$ 却这样表示。
       \begin{equation*}
         \begin{mat}
           0  &1  \\
           0  &0
         \end{mat}
       \end{equation*}
     \end{answer}
   \recommended \item
     证明若一个变换的值域空间与零空间相同，
     则其定义域的维数是偶数。
     \begin{answer}
       设 $\map{t}{V}{V}$ 是该变换。
       若 \( \rank (t)=\nullity (t) \)，则等式
       \( \rank(t)+\nullity(t)=\dim (V) \) 表明
       \( \dim (V) \) 是偶数。
     \end{answer}
   \item
     证明若两个幂零矩阵可交换，则它们的乘积与
     和也是幂零的。
     \begin{answer}
       矩阵要是幂零的，它们必须是方阵。
       要可交换，它们必须大小相同。
       于是它们的乘积与和都有定义。

       把这两个矩阵称为 \( A \) 与 \( B \)。
       要看出 \( AB \) 是幂零的，计算
       $
          (AB)^2=ABAB=AABB=A^2B^2$，
       以及
          $(AB)^3=A^3B^3$，依此类推，
       又因 \( A \) 是幂零的，该乘积最终为零。

       和的情形类似；利用二项式定理。
      \end{answer}
   % 2014-Dec-19 These two questions and answers are fine; I commented them to make the pages fit better.
   % \item
   %   Consider the transformation of \( \matspace_{\nbyn{n}} \) given
   %   by \( t_S(T)=ST-TS \) where \( S \) is an \( \nbyn{n} \) matrix.
   %   Prove that if \( S \) is nilpotent then so is \( t_S \).
   %   \begin{answer}
   %     Some experimentation gives the idea for the proof.
   %     Expansion of the second power
   %     \begin{equation*}
   %       t^2_S(T)=S(ST-TS)-(ST-TS)S=S^2-2STS+TS^2
   %     \end{equation*}
   %     the third power
   %     \begin{align*}
   %       t^3_S(T)
   %         &=S(S^2-2STS+TS^2)-(S^2-2STS+TS^2)S  \\
   %         &=S^3T-3S^2TS+3STS^2-TS^3
   %     \end{align*}
   %     and the fourth power
   %     \begin{align*}
   %       t^4_S(T)
   %         &=S(S^3T-3S^2TS+3STS^2-TS^3)-(S^3T-3S^2TS+3STS^2-TS^3)S  \\
   %         &=S^4T-4S^3TS+6S^2TS^2-4STS^3+TS^4
   %     \end{align*}
   %     suggest that the expansions follow the Binomial Theorem.
   %     Verifying this by induction on the power of $t_S$ is routine.
   %     This answers the question because, where the index of nilpotency of
   %     $S$ is $k$, in the expansion of $t^{2k}_S$
   %     \begin{equation*}
   %       t^{2k}_S(T)=\sum_{0\leq i\leq 2k}(-1)^i\binom{2k}{i} S^iTS^{2k-i}
   %     \end{equation*}
   %     for any $i$ at least one of the $S^i$ and $S^{2k-i}$
   %     has a power higher than $k$, and so the term gives the zero matrix.
   %   \end{answer}
   % \item
   %   Show that if \( N \) is nilpotent then \( I-N \) is
   %   invertible.
   %   Is that `only if' also?
   %   \begin{answer}
   %     Use the geometric series:
   %     $
   %        I-N^{k+1}=(I-N)(N^k+N^{k-1}+\cdots+I)
   %     $.
   %     If \( N^{k+1} \) is the zero matrix then we have a right inverse for
   %     \( I-N \).
   %     It is also a left inverse.

   %     This statement is not `only if' since
   %     \begin{equation*}
   %        \begin{mat}[r]
   %           1  &0  \\
   %           0  &1
   %        \end{mat}
   %        -\begin{mat}[r]
   %           -1  &0  \\
   %           0  &-1
   %        \end{mat}
   %     \end{equation*}
   %     is invertible.
   %    \end{answer}
 %  \recommended \item
 %    Show that when a nilpotent matrix \( T \) is in canonical form
 %    then the number of
 %    blocks of \( \nbyn{(k+1)} \) matrices that are all zeros
 %    except for subdiagonal ones is
 %    \( 2\,(\nullity (T^k))
 %      -\nullity (T^{k+1})-\nullity (T^{k-1}) \).
\index{nilpotent|)}
\end{exercises}
